\scanpage{71}
% archive sheet number: 70
% L1-3
Regardons maintenant l'\uline{hom. norme}, et soit $V$ son noyau
\[
e\longrightarrow V\longrightarrow(G_m)_1\mathbin{-}G_m\longrightarrow e
\]
% L4-8
\uline{Mais}, $V$ \uline{n'est pas lisse sur $k$}, \deleted{\unsure{le}\ \illegible{1}\ \unsure{groupe}} \deleted{Mais [dans le cas où \unsure{donc}} pas plus que $V_{\mathrm{réd}}$ [car autrement ce serait le radical \inserted{\uline{unipotent}} de $G_m$].
\margin{dans le cas précédent envisagé où $H^1(k,U)\ne0$}
D'ailleurs, dans le cas où $[k':k]=p$, on voit facilement que $V$ est réduit. D'ailleurs, dans le cas où $k'^p\subset k$, on voit facilement que $V$ \uline{est aussi le noyau de la puissance $p$-ième dans} $G_{m\,1}$, donc on a un isom
% slanted
\margin{\uuline{\unsure{préciser}} \unsure{prouver par descente} ---}
\[
V\simeq\prod_{S'/S}(\unsure{\boldsymbol{\mu}|\mu}_p)_{S'}
\]
% L13-17
\deleted{Dans tous cas}, ce qui prouve qu'on a en fait une suite exacte
\[
0\mathbin{-}\boldsymbol{\mu}_p\mathbin{-}V\longrightarrow U\longrightarrow0
\]
(extension non triviale ---).

% L18-22
Notons maintenant que si $G'$ est un groupe alg. \inserted{\unsure{connexe simple|connexe réduit}} sur $k'$ \deleted{\unsure{un groupe}} ayant un centre isom. à $(\boldsymbol{\mu}_p)_{k'}$
\margin{[p.ex. $\operatorname{SL}(p)\deleted{\unsure{(k')}}_{k'}$]}
alors le centre de $G_1$ est
\[
Z_1=V
\]
donc $(\bar Z_1)_{\mathrm{réd}}$ \uline{n'est pas définissable sur} $k$
